% Figure 2: Type System Overview
% Hierarchical type system for OTLP with key typing rules

\begin{figure}[t]
\centering
\begin{tikzpicture}[
    node distance=1.2cm and 1.5cm,
    type/.style={rectangle, draw=black!70, fill=blue!10, thick, minimum width=2.8cm, minimum height=0.9cm, rounded corners, font=\small},
    rule/.style={rectangle, draw=black!60, fill=gray!10, thick, minimum width=5.5cm, minimum height=0.7cm, rounded corners, font=\footnotesize},
    arrow/.style={->, >=stealth, thick}
]

% Type hierarchy
\node[type] (trace) at (0,0) {\textbf{Trace}\\$\langle$tid, spans$\rangle$};
\node[type, below left=of trace] (span) {Span\\$\langle$sid, attrs$\rangle$};
\node[type, below right=of trace] (resource) {Resource\\$\langle$attrs$\rangle$};
\node[type, below=1.5cm of span] (attr) {Attribute\\$\langle$k, v, type$\rangle$};

% Relationships
\draw[arrow] (trace) -- node[left, font=\tiny] {contains*} (span);
\draw[arrow] (trace) -- node[right, font=\tiny] {has} (resource);
\draw[arrow] (span) -- node[left, font=\tiny] {has*} (attr);
\draw[arrow] (resource) -- node[right, font=\tiny] {has*} (attr);

% Key typing rules (on the right)
\node[rule, right=1cm of trace] (rule1) {
    \textbf{T-Trace:} $\Gamma \vdash$ trace well-typed\\
    \tiny if tid unique $\land$ spans valid
};

\node[rule, below=0.3cm of rule1] (rule2) {
    \textbf{T-Span:} $\Gamma \vdash$ span well-typed\\
    \tiny if sid unique $\land$ parent exists
};

\node[rule, below=0.3cm of rule2] (rule3) {
    \textbf{T-Temporal:} $t_{\text{start}} < t_{\text{end}}$\\
    \tiny parent.start $\leq$ child.start
};

\node[rule, below=0.3cm of rule3] (rule4) {
    \textbf{T-Context:} context propagates\\
    \tiny if span has parent $\Rightarrow$ inherits ctx
};

% Soundness theorem box
\node[draw=red!70, fill=red!10, thick, text width=5.3cm, rounded corners, font=\small, below=0.5cm of rule4] (soundness) {
    \textbf{Theorem (Type Soundness):}\\
    If $\Gamma \vdash t : \text{Trace}$, then\\
    \textbf{Progress:} $t$ is value or $t \rightarrow t'$\\
    \textbf{Preservation:} $t \rightarrow t' \Rightarrow \Gamma \vdash t' : \text{Trace}$
};

% Properties box (bottom)
\node[draw=blue!70, fill=blue!5, thick, text width=11cm, rounded corners, font=\footnotesize, below=1.8cm of attr] (properties) {
    \textbf{Key Properties Enforced by Type System:}
    \begin{itemize}[leftmargin=*, itemsep=0pt]
        \item \textbf{Uniqueness:} trace IDs and span IDs are unique within scope
        \item \textbf{Hierarchy:} parent-child relationships form a valid tree
        \item \textbf{Temporality:} timestamps maintain causal order
        \item \textbf{Completeness:} required attributes present, contexts propagated
        \item \textbf{Type Safety:} attribute values match declared types
    \end{itemize}
};

\end{tikzpicture}
\caption{\textbf{OTLP Type System Overview.} The type hierarchy (left) shows the core OTLP data types with containment relationships. Key typing rules (right) ensure correctness properties including ID uniqueness, valid hierarchies, temporal ordering, and context propagation. The type soundness theorem guarantees that well-typed OTLP traces maintain their properties during execution (formalized in 1,500 lines of Coq).}
\label{fig:typesystem}
\end{figure}

